% Part: computability
% Chapter: recursive-functions
% Section: halting-problem

\documentclass[../../../include/open-logic-section]{subfiles}

\begin{document}

\olfileid{cmp}{rec}{hlt}
\olsection{مسئلهٔ توقف}

\emph{مسئلهٔ توقف} در صورت کلی مسئلهٔ تصمیم‌گیری دربارهٔ این است که
با داشتن مشخصهٔ~$e$ (مثلاً برنامه) از یک تابع محاسبه‌پذیر و عدد~$n$،
آیا محاسبهٔ تابع بر ورودی~$n$ متوقف می‌شود، یعنی نتیجه‌ای تولید
می‌کند، یا نه. آلن تورینگ در نتیجه‌ای مشهور ثابت کرد که خود این مسئله
را نمی‌توان با تابعی محاسبه‌پذیر حل کرد؛ یعنی تابع
\[
h(e, n) =
\begin{cases}
1 & \text{اگر محاسبهٔ $e$ روی ورودی $n$ متوقف شود}\\
0 & \text{در غیر این صورت،}
\end{cases}
\]
محاسبه‌پذیر نیست.

در زمینهٔ توابع بازگشتی جزئی، نمایهٔ~$e$ که در قضیهٔ صورت نرمال
کلینی داده می‌شود می‌تواند نقش مشخصهٔ برنامه را بازی کند. اگر $f$
تابعی بازگشتی جزئی باشد، هر~$e$ای که معادلهٔ قضیهٔ صورت نرمال برایش
برقرار است نمایه‌ای از~$f$ است. برای عدد داده‌شدهٔ~$e$، قضیهٔ صورت
نرمال می‌گوید
\[
\cfind{e}(x) \simeq U(\mu s \; T(e, x, s))
\]
بازگشتی جزئی است، و برای هر تابع بازگشتی جزئی
$f\colon\Nat\to\Nat$ عددی $e\in\Nat$ وجود دارد که به ازای همهٔ
$x\in\Nat$، $\cfind{e}(x)\simeq f(x)$. در واقع برای هر چنین~$f$ای
نه‌فقط یک، بلکه بی‌نهایت~$e$ از این نوع وجود دارد. \emph{تابع
توقف}~$h$ چنین تعریف می‌شود:
\[
h(e, x) =
\begin{cases}
1 & \text{اگر $\cfind{e}(x) \fdefined$}\\
0 & \text{در غیر این صورت.}
\end{cases}
\]
توجه کنید که اگر $\cfind{e}(x)\fundefined$، آنگاه $h(e,x)=0$؛ و نیز
اگر~$e$ اصلاً نمایهٔ یک تابع بازگشتی جزئی نباشد همین مقدار را دارد.

\begin{thm}
\ollabel{thm:halting-problem}
تابع توقف~$h$ بازگشتی جزئی نیست.
\end{thm}

\begin{proof}
اگر $h$ بازگشتی جزئی بود، می‌توانستیم تعریف کنیم:
\[
d(y) =
\begin{cases}
1 & \text{اگر $h(y, y) = 0$}\\
\umin{x}{x \neq x} & \text{در غیر این صورت.}
\end{cases}
\]
چون هیچ عدد $x$ای در $x\neq x$ صدق نمی‌کند،
$\umin{x}{x\neq x}$ وجود ندارد؛ پس $d(y)\fundefined$ اگر و تنها اگر
$h(y,y)\neq0$. از این تعریف نتیجه می‌شود:
\begin{enumerate}
\item $d(y)\fdefined$ اگر و تنها اگر $\cfind{y}(y)\fundefined$، یا
  $y$ نمایهٔ یک تابع بازگشتی جزئی نباشد.
\item $d(y)\fundefined$ اگر و تنها اگر $\cfind{y}(y)\fdefined$.
\end{enumerate}
اگر $h$ بازگشتی جزئی بود، $d$ نیز بازگشتی جزئی می‌بود. پس بنا بر
قضیهٔ صورت نرمال کلینی نمایه‌ای چون~$e_d$ می‌داشت. مقدار
$h(e_d,e_d)$ را در نظر بگیرید. دو حالت ممکن، $0$ و~$1$، وجود دارد:
\begin{enumerate}
\item اگر $h(e_d,e_d)=1$، آنگاه $\cfind{e_d}(e_d)\fdefined$. اما
  $\cfind{e_d}\simeq d$ و $d(e_d)$ تعریف شده است اگر و تنها اگر
  $h(e_d,e_d)=0$. پس $h(e_d,e_d)\neq1$.
\item اگر $h(e_d,e_d)=0$، آنگاه یا $e_d$ نمایهٔ تابعی بازگشتی جزئی
  نیست، یا چنین نمایه‌ای است و $\cfind{e_d}(e_d)\fundefined$. اما
  باز هم $\cfind{e_d}\simeq d$ و $d(e_d)$ تعریف‌نشده است اگر و تنها
  اگر $\cfind{e_d}(e_d)\fdefined$.
\end{enumerate}
نتیجه این است که $e_d$ برخلاف فرض نمی‌تواند نمایهٔ یک تابع بازگشتی
جزئی باشد. اما اگر $h$ بازگشتی جزئی بود، $d$ نیز چنین می‌بود و تعریف
$e_d$ به‌عنوان نمایهٔ آن مجاز می‌شد. باید نتیجه بگیریم که $h$
نمی‌تواند بازگشتی جزئی باشد.
\end{proof}

\end{document}
